\section{Weighted blocks beyond cacti}
\label{sec:a-weighted-blocks}

The small-face theorem in \cref{thm:a-wcac-faces} applies to every fixed
maximum block rank. Its remaining algorithmic issue is summation: many
blocks can depend on the same few unfixed nominations. This section
resolves that issue for one shared nomination parameter, including
arbitrarily many noncactus blocks. It also permits any fixed number of
shared nomination parameters when the total rank of the noncactus blocks
is fixed. Both results include independent coefficient intervals and
rational near-optimal original scenarios. The unrestricted combination
of several shared parameters and arbitrarily many higher-rank blocks is
delimited at the end.

Except in \cref{cor:a-wblk-line-polynomial}, laws in this section are
quadratic and the allowed coefficient
models are exactly those in \cref{thm:a-wcac-affine}: fixed positive
asymmetric coefficients, independent symmetric coefficient intervals, or
independent directional coefficient intervals. No operating constraints
filter the scenarios. Put $W=c^\top\pi$, where $c$ is rational and balanced.
All complexity statements refer to the Turing model and accuracy bits.

\subsection{Exact local value graphs}

\begin{lemma}[Conditional block values]
\label{lem:a-wblk-local}
Fix $d,r_0$. Let $b(z)=b^0+Hz$ be balanced on a nonempty compact rational
polytope $P\subseteq\R^d$, and suppose every block has rank at most $r_0$.
One can construct, for every block $B$, a rational quantifier-free formula
$\Gamma_B(z,v)$ for the graph of its maximum weighted contribution
$V_B(z)$, conditional on $z$ and optimized over its allowed coefficients.
Formula sizes and coefficient bit lengths are polynomial, and degrees
are bounded in terms of $d,r_0$. Moreover,
\[
 \max_{z,\theta}W=\max_{z\in P}\sum_B V_B(z).
\]
Each $V_B$ has a computable rational absolute bound and a Lipschitz
constant of polynomial bit length. At a rational $z\in P$, its value and
an optimizing local coefficient profile can be computed as real algebraic
data of polynomial representation size, without joining fields from
different blocks.
\end{lemma}
\begin{proof}
Choose a spanning tree formed by a spanning tree in each block, and a
rational tree flow $w$ with $Aw=c$. Then
$W=\sum_e w_eg_e(x_e)$ by \eqref{eq:a-wcac-weights}. A block's effective
nomination $b^B(z)$ is obtained by summing nominations over the components
attached to each of its vertices. These components partition $V$, so
$b^B$ is rational affine and balanced. Its physical state depends only
on $b^B$ and its own coefficients. Block independence therefore proves
the displayed identity, with each local contribution
$\sum_{e\in B}w_eg_e(x_e)$.

Tree routing and a fundamental-cycle basis give
$x_e=X_e(z,q_B)$, with $r_B\le r_0$ circulation coordinates. A rational
bound $B_0\ge1$ on total positive nominations bounds every physical flow;
because the tree routing vanishes on chords, $|q_{B,i}|\le B_0$.
Enumerate the realizable flow-sign cells in $(z,q_B)$ and use their
closures. On such a cell put $Q_e=s_eX_e^2$. At fixed coefficients the
cycle equations and local value are quadratic. For uncertain coefficients
apply \cref{lem:a-blk-boxlp}, with the $r_B$ cycle equations as linking
rows and $w_eQ_e$ as the single objective row. In the directional model
use the interval of the selected sign; unused directional coordinates
can take arbitrary allowed values. At zero flow both directional terms
vanish, so closed sign-cell overlaps are exact. The symmetric model
retains one original coefficient per edge.

This gives a polynomial-size fixed-degree formula for attainable
$(z,q_B,v)$. Eliminate $q_B$ to obtain $\Psi_B(z,v)$. The desired graph is
\[
 \Psi_B(z,v)\ \wedge\
 \neg\exists v'\bigl(\Psi_B(z,v')\wedge v'>v\bigr).
\]
All projections have a fixed number of variables. The bounds in
\cref{subsec:a-pre-computation} prove the representation claims, including
singular states and equality cases. Every fiber is nonempty and compact:
all coefficients are in compact boxes, and the unique physical state
depends continuously on them. At rational $z$, the same computation and
\cref{lem:a-blk-boxlp} recover an exact local optimizer. A bridge needs no
circulation variable and is included in this argument.

For the continuity bounds, let $H^B$ be the linear part of $b^B$ and put
\[
 T_B=\sum_{v,j}|H^B_{vj}|,\qquad
 K_B=\beta_U B_0T_B\sum_{e\in B}|w_e|,\qquad
 M_B=\beta_U B_0^2\sum_{e\in B}|w_e|.
\]
Here $\beta_U$ bounds every allowed directional coefficient. At fixed
coefficients, \eqref{eq:a-wcac-flow-nomination}, applied inside the block,
and the edge-law Lipschitz constant $2\beta_U B_0$ give
\[
 |V_B(z)-V_B(z')|\le K_B\|z-z'\|_\infty,
 \qquad |V_B(z)|\le M_B.
\]
For the first inequality in the uncertain case, apply the fixed-profile
inequality at a maximizer for $z$, compare that same allowed profile at
$z'$, and reverse the roles. Thus taking a maximum preserves the bound.
All displayed constants have polynomial bit length, using a rational
bounding box for $P$. They do not require lower bounds on nonzero flows
or on derivatives.
\end{proof}

\subsection{A rational approximation lemma for one parameter}

The next lemma concerns a scalar semialgebraic function, rather than a
special radical formula. A quantifier-free graph description means an
explicit Boolean formula in rational polynomial sign tests whose
solutions over the stated interval are exactly $(t,f(t))$. Degrees may
grow when the polynomials are densely encoded.

\begin{lemma}[Uniform rational approximation on an interval]
\label{lem:a-wblk-univariate}
Let $f:[0,1]\to\R$ have a rational quantifier-free semialgebraic graph
description, and suppose rational $K,M_0\ge1$ satisfy
$|f(t)-f(t')|\le K|t-t'|$ and $|f(t)|\le M_0$.
Given rational $0<\eta\le1$, one can construct rational closed intervals
covering $[0,1]$ and a rational polynomial on each interval whose error
from $f$ is at most $\eta$ throughout that interval. The number of
intervals, polynomial degrees, dense coefficient encodings, and
construction time are polynomial in the graph description, the encodings
of $K,M_0$, and the accuracy bits. Adjacent polynomials need not agree at
their shared endpoints.
\end{lemma}
\begin{proof}
We give all approximation and precision bounds. Polynomial arithmetic,
subresultants, root isolation, and fixed-dimensional real elimination are
the classical algorithms described in
\citet{basu2006-algorithms-in-real-algebraic-geometry} and
\cref{subsec:a-pre-computation}.

\emph{An annihilating polynomial and its exceptional parameters.}
Discard identically zero polynomials from the graph formula, replacing
their tests by their truth values. Multiply its polynomials that depend
on the value variable $y$, remove the content in $\Q[t]$, and take the
squarefree part in $\Q(t)[y]$. Clear denominators and remove content
again to obtain a nonzero primitive $F(t,y)\in\Q[t,y]$ with positive
$y$-degree. This polynomial satisfies
\begin{equation}
 F(t,f(t))=0\quad(0\le t\le1).
 \label{eq:a-wblk-annihilator}
\end{equation}
Indeed, outside the finitely many roots of the nonzero $t$-only
polynomials and the contents, some value-dependent polynomial must
vanish at $(t,f(t))$. Otherwise all polynomial signs are locally constant,
and the graph formula would contain a two-dimensional neighborhood,
contrary to being a function graph. Taking a primitive squarefree part
preserves these zeros away from finitely many additional specializations.
Continuity extends the resulting identity to every remaining parameter.
This also proves that there was a value-dependent factor to start with.

The product has polynomial total degree and coefficient length. Gcd and
squarefree computation over $\Q(t)$ have the same polynomial bounds:
subresultant coefficients are determinants of polynomial-size matrices
of univariate polynomials, with polynomial degree and bit length. Taking
primitive parts and exact quotients preserves polynomial bounds. Thus
the following construction does not require a bounded degree for $F$.

Write $F(t,y)=\sum_{j=0}^D a_j(t)y^j$, with $a_D\ne0$. The polynomial
\[
 H_*(t)=a_D(t)\operatorname{Res}_y(F,\partial_yF)
\]
is nonzero because $F$ is squarefree over $\Q(t)$. Its degree $E$ and
coefficient lengths are polynomial. Its zeros include every finite
parameter where a root can collide or the leading coefficient can vanish.
Compute the distinct real parts of its complex roots. For example,
eliminate $v$ from
$\Re H_*(u+iv)=\Im H_*(u+iv)=0$ and isolate the resulting real $u$.
This has two real variables and polynomial degree, so takes polynomial
bit time. There are at most $E$ such real parts; nonreal roots must be
included. When $H_*$ is constant the list is empty.

\emph{Short intervals around exceptional real parts.}
Choose $\rho=2^{-s}$, where $s$ is the least integer $s\ge0$ such that
\[
 \rho\le\min\{1/16,\eta/[64K(E+1)]\}
\]
holds. If $A$ denotes the displayed upper bound, then $A/2<\rho\le A$,
so $\rho$ has polynomial bit length.
Keep every exceptional real part in $[-1,2]$,
enclose it in a rational interval $[l,u]$ of width at most $\rho$, and
pad that interval to $[l-\rho,u+\rho]$. Include $[-\rho,\rho]$ and
$[1-\rho,1+\rho]$, intersect the union with $[0,1]$, and merge overlaps.
Call this union $J$. Its total length is at most
$3E\rho+4\rho\le4(E+1)\rho\le\eta/(16K)$.
On each component use a rational approximation, within $\eta/8$, to
$f$ at the rational midpoint. The Lipschitz bound makes its error
throughout that component less than $\eta/4$. Values at rational
arguments are computed by the graph formula and real-algebraic sampling,
then isolated to the required rational accuracy. These operations have
polynomial cost in the graph data and requested precision.

\emph{Analytic disks and rational panels.}
For a complementary gap $[a,b]$ put
$d(t)=\min\{t-a+\rho,b-t+\rho\}$.
Every exceptional real part is at distance at least $d(t)$ from $t$.
For retained real parts this follows from the padding, also after merging;
omitted ones are outside $[-1,2]$, whereas $d(t)\le1/2+\rho<1$.
Split the gap at its midpoint. Starting at $x=a$, make left-half panels
\[
 [x,y],\qquad y=\min\{(a+b)/2,x+(x-a+\rho)/32\},
\]
and reflect this rule for the right half. Before the last panel,
$x-a+\rho$ grows by $33/32$. Hence there are
$O(1+\log(1/\rho))$ panels per gap, all with polynomial rational encoding.
If $\tau$ and $h$ are a panel's midpoint and half-width, then
\[
 h\le d(\tau)/64,
 \qquad R=d(\tau)/4,
 \qquad h\le R/16.
\]
The closed complex disk $|t-\tau|\le R$ stays at distance at least
$3d(\tau)/4\ge3\rho/4$ from every zero of $H_*$.

All roots of $F(t,\cdot)$ are simple on a neighborhood of this disk,
and $a_D(t)$ is nonzero. They form a finite unramified covering there:
local implicit branches exist, and a uniform polynomial root bound
prevents escape to infinity along a compact path. Since a disk is simply
connected, the branch through $(\tau,f(\tau))$ is single valued and
holomorphic on a neighborhood of the closed disk. On the real panel it
equals $f$, because a continuous choice of distinct roots cannot change
branches. This step allows nonsmooth branch switches inside $J$ and
does not assume that the original graph is irreducible.

Here is a uniform bound on every such holomorphic branch. Write $k_0$ for
the degree of $a_D$ and $\lambda\ne0$ for its leading coefficient. All
its roots are zeros of $H_*$. Thus throughout these disks
\[
 |a_D(t)|\ge |\lambda|(3\rho/4)^{k_0}.
\]
Also $|t|<2$. Expanding the $a_j$ gives a rational $T\ge1$ bounding
$|a_j(t)|$ there. The elementary Cauchy polynomial-root bound therefore
gives
\begin{equation}
 |f_{\rm hol}(t)|\le
 M:=1+\frac{T}{|\lambda|(3\rho/4)^{k_0}}.
 \label{eq:a-wblk-complex-bound}
\end{equation}
The logarithm and rational encoding length of $M$ are polynomial, even
if a complex singularity is extremely close to the real interval.

\emph{Rational interpolation with a certified error.}
Choose the least integer $q\ge1$ with $M2^{-q}\le\eta/8$; its numerical
value is polynomially bounded. On a panel take the $q+1$ equally spaced rational nodes
$t_j=\tau-h+2hj/q$. Let $I_q$ interpolate the exact values $f(t_j)$.
The contour interpolation remainder on $|t-\tau|=R$ yields, for real
$x$ in the panel,
\begin{align*}
 |f(x)-I_q(x)|
 &\le\frac{MR}{R-h}
       \left(\frac{2h}{R-h}\right)^{q+1}\\
 &\le\frac{16M}{15}(2/15)^{q+1}
 \le M2^{-q}\le\eta/8.
\end{align*}
For completeness, the remainder is the contour integral of
$f_{\rm hol}(t)\omega(x)/[(t-x)\omega(t)]$ divided by $2\pi i$, where
$\omega(t)=\prod_j(t-t_j)$. Bound each numerator factor by $2h$, each
denominator factor by $R-h$, and the contour length by $2\pi R$.
At a node the error is zero directly.

No algebraic coefficient field is needed to construct the approximation.
Replace each $f(t_j)$ by a rational value with error at most
\[
 \delta=\frac{\eta}{8(2q)^q}.
\]
For the Lagrange basis polynomials $L_j$ on these nodes,
\[
 \sum_{j=0}^q |L_j(x)|
 \le\sum_{j=0}^q\frac{q^q}{j!(q-j)!}
 =\frac{(2q)^q}{q!}\le(2q)^q.
\]
The rational interpolant consequently differs from $I_q$ by at most
$\eta/8$, and from $f$ by at most $\eta/4$. Its required node precision
has $O(\log(1/\eta)+q\log(q+1))$ bits. Root isolation computes the node
values independently. Expanding the rational Lagrange formula requires
polynomially many arithmetic operations with polynomial-bit numbers:
each product has $q$ factors, each node difference has polynomial
encoding, and summing $q+1$ such polynomials preserves that bound.
Thus the dense rational coefficient representation also has polynomial
size. The constant pieces on $J$ and these analytic panels cover the
whole interval and prove every assertion.
\end{proof}

\begin{corollary}[Sums of scalar semialgebraic responses]
\label{cor:a-wblk-scalar-sum}
For a list of functions satisfying \cref{lem:a-wblk-univariate}, both
extrema of their sum on $[0,1]$ admit rational enclosing intervals of
width at most $\eps$ and rational $\eps$-optimal arguments in time
polynomial in the combined input and accuracy bits. The number of
summands and their algebraic degrees may grow.
\end{corollary}
\begin{proof}
For $m$ summands apply the lemma with $\eta=\eps/(16\max\{1,m\})$,
after replacing $\eps$ by $\min\{\eps,1\}$. Sort all rational panel
endpoints. On each common closed subinterval select one valid polynomial
per summand and add them. This gives a rational polynomial with uniform
error at most $a=\eps/16$ from the true sum. The number of subintervals,
degrees, and coefficient lengths are polynomial; endpoints are covered
by either adjoining valid choice.

Optimize each rational polynomial on its closed interval by univariate
real-algebraic methods. These compute a rational interval $[L,U]$ of
width at most $\eps/8$ containing the largest surrogate maximum and an
algebraic argument whose surrogate value is at least $L-\eps/16$.
Then $[L-a,U+a]$ encloses the true maximum and has width at most
$\eps/4$. The sample's true value is at least $L-\eps/16-a$, so its
loss is at most $5\eps/16$. Isolate the sample in an interval of width
$\eps/[4\sum_i K_i]$ and choose a rational point in its intersection
with $[0,1]$. The true sum loses at most $\eps/4$ by Lipschitz
continuity. Every isolation precision has polynomial bit length. A
constant sum can instead use any rational argument; the empty sum is
zero. Negating every function proves the minimum statement.
\end{proof}

\subsection{One shared nomination parameter}

\begin{theorem}[Weighted bounded blocks on an affine nomination line]
\label{thm:a-wblk-line}
Fix $r_0$. On a connected graph of maximum block rank at most $r_0$,
let $b(t)=b^0+ht$ be a rational balanced nomination family on a nonempty
compact rational interval. Objective support, the number of noncactus
blocks, and total cycle rank are unrestricted. In all three coefficient
models stated above, each joint extremum of $W$ admits a rational
enclosing interval of width at most $\eps$ and a rational $\eps$-optimal
original parameter scenario in accuracy-bit polynomial time.
\end{theorem}
\begin{proof}
A singleton interval reduces to independent exact local computations
and separately refined value intervals. Otherwise a rational affine
change of variable makes the interval $[0,1]$; it preserves polynomial
encoding. Apply \cref{lem:a-wblk-local} with $d=1$ and then
\cref{cor:a-wblk-scalar-sum}, using tolerance $\eps/2$. Increase any zero
local bound to one to meet the approximation lemma's conventions. This
gives a certified interval and a rational parameter $t'$ whose conditional
joint value loses at most $\eps/2$.

At this rational $t'$, recover a coefficient maximizer separately in
each block by \cref{lem:a-wblk-local}. These local choices form a global
physical scenario because their nominations agree at the attachments
and their cycle equations hold in every block. The output of this
intermediate step may be algebraic. If $k\le2m$ original coefficients
are uncertain, isolate each to width $\delta$, intersect with its own
original rational interval, and choose a rational value there. Choose
\[
 \|c\|_1 B_0^2 k\delta\le\eps/4.
\]
If its left coefficient is zero there is no loss to control. The
all-graph estimate \eqref{eq:a-wcac-joint-continuity} bounds the objective
loss by $\eps/4$. Shared symmetric coefficients are rounded once;
directional coordinates are rounded separately, including inactive ones.
Fixed and singleton coefficients are retained exactly. Isolation and
rounding take polynomial time because local representations have bounded
degree and polynomial length. There is no common field containing all
block roots. The resulting rational nomination and coefficients satisfy
their original unfiltered domains exactly, and their physical state has
the required objective. Negate $c$ for the minimum.
\end{proof}

The theorem resolves the univariate summation issue even when the local
functions are not expressed by radicals. It does not turn the exact
comparison of their sum into a polynomial-time problem. Certified
surrogates, not separation bounds for exact sums, are used throughout.

\begin{corollary}[Fixed dense piecewise-polynomial laws]
\label{cor:a-wblk-line-polynomial}
The fixed-law conclusion of \cref{thm:a-wblk-line} remains true for the
continuous, strictly increasing, densely encoded rational
piecewise-polynomial edge laws vanishing at zero from
\cref{subsec:a-law-model}. Degrees and piece counts may grow with the
input. This extension asserts fixed laws, with rational nomination output.
\end{corollary}
\begin{proof}
Use the same affine flow coordinates $(t,q_B)$ and polynomial-bit bound
$B_0\ge1$. Partition by every affine equation $X_e(t,q_B)=a$, where $a$
is an edge-law breakpoint. There are polynomially many realizable cells
in fixed dimension $1+r_0$. On each closed cell the cycle equations and
weighted contribution are rational polynomials whose degrees and dense
encodings are polynomial in the input. Continuity makes the formulas
agree at breakpoints. Conservation and cycle consistency still characterize
the unique physical state by \cref{thm:a-pre-existence}. Fixed-dimensional
elimination therefore constructs a polynomial-size graph for each local
value, now of polynomial rather than fixed degree.

Absolute coefficient sums over all edge pieces give polynomial-bit
rational bounds $G_0$ and $L_0$ on their values and one-sided derivatives
on $[-B_0,B_0]$. For a piece $\sum_j a_jx^j$ one can use
$\sum_j|a_j|B_0^j$ and $\sum_{j\ge1}j|a_j|B_0^{j-1}$, respectively.
Continuity across finitely many breakpoints makes $L_0$ a Lipschitz
constant for the complete edge law on that interval. The acyclic
difference-flow proof of \eqref{eq:a-wcac-flow-nomination} uses only
strict monotonicity, so remains valid. Consequently each block's value
has absolute bound $G_0\sum_{e\in B}|w_e|$ and nomination-parameter
Lipschitz constant $(L_0/2)T_B\sum_{e\in B}|w_e|$.
\Cref{lem:a-wblk-univariate,cor:a-wblk-scalar-sum} allow growing dense
degree and supply the claimed value interval and rational argument.
No law approximation or coefficient rounding is needed. Sparse
binary-encoded large exponents are outside this statement.
\end{proof}

\subsection{Several parameters with finitely much noncactus rank}

Define
\[
 \kappa(G)=\sum_{B:\,r_B\ge2}r_B.
\]
Thus $\kappa(G)=0$ on cacti. A graph can have fixed $\kappa(G)>0$ and
arbitrarily large total rank, for example a subdivided theta block with
arbitrarily many attached cycle blocks. Neither the lengths of the
exceptional blocks nor the number of their uncertain coefficients are
bounded by $\kappa(G)$.

\begin{theorem}[A fixed total rank in the noncactus blocks]
\label{thm:a-wblk-hybrid}
Fix $d,\kappa_0$. Suppose $\kappa(G)\le\kappa_0$ and nominations form a
rational affine balanced family $b(z)=b^0+Hz$ on a nonempty compact
rational polytope $P\subseteq\R^d$. In all three coefficient models,
with unrestricted objective support, both joint extrema of $W$ admit
rational intervals of width at most $\eps$ and rational $\eps$-optimal
original parameter scenarios in accuracy-bit polynomial time.
\end{theorem}
\begin{proof}
Call the blocks of rank at least two exceptional. Choose the blockwise
spanning tree and objective weights $w$ as in \cref{lem:a-wblk-local}.
Retain the circulations $q$ in all exceptional blocks; there are
$\kappa=\kappa(G)$ of them. Every exceptional-edge flow $X_e(z,q)$ is
rational affine, and $|q_i|\le B_0$. Thus the exact exceptional core has
fixed dimension $d+\kappa$, even if it contains arbitrarily many edges.

Enumerate the realizable exceptional flow-sign cells in these variables
and use closed cells intersected with $P\times[-B_0,B_0]^\kappa$.
For each cell use \cref{lem:a-blk-boxlp} with all $\kappa$ exceptional
cycle equations as linking rows and the \emph{one} value row
\begin{equation}
 v=\sum_{e\text{ exceptional}}w_e\beta_e s_eX_e(z,q)^2.
 \label{eq:a-wblk-exceptional-row}
\end{equation}
The linking matrix is block diagonal in its cycle rows. The box lemma
eliminates every exceptional coefficient and gives a quantifier-free
formula $\Psi(z,q,v)$ for exactly the attainable cores and total
exceptional values. Its degree is bounded in terms of $d,\kappa_0$,
and its size and coefficient lengths are polynomial. Arbitrary objective
support changes $w$, but does not add linking or value coordinates.
Directional intervals are selected by the flow sign as in
\cref{lem:a-wblk-local}; fixed coefficients are singleton intervals.
If there are no exceptional blocks, simply take $v=0$.

The remaining blocks are cycles and bridges. Their through-nominations
depend on $z$, and not on the exceptional circulations or coefficients:
block aggregation determines them from the original nominations alone.
Use the local candidate construction and rational approximations of
\cref{sec:a-weighted-cactus} for these blocks. That construction only
uses their affine local nominations, their one cycle equation, and their
weights $w_e$; it does not require the other blocks of the graph to be
cycles. In fixed dimension $d$ it gives polynomially many common sign
parts with rational surrogates $S(z)$ for their summed conditional maximum.
Allocate local errors so that both the true conditional maximum and a
recoverable selected local scenario differ from $S(z)$ by at most
$a=\eps/16$. These are the two guarantees proved after
\cref{lem:a-wcac-candidates}. Bridge endpoint choices are exact.

On a common part optimize
\[
 v+S(z)\quad\text{subject to}\quad\Psi(z,q,v).
\]
The number of variables is $d+\kappa+1$. Although the surrogate's
degree grows with requested precision, its degree and dense encoding
are polynomial. Denominators have known nonzero signs on each part;
clear them exactly as in \cref{sec:a-weighted-cactus}. The box lemma
was used \emph{before} this approximation, with degree-two coefficient
data, so its fixed-degree hypothesis has not been changed.

The projected exceptional set is compact. Individual surrogate sign
parts can be open, but the physical objective is bounded by
\eqref{eq:a-wcac-value-bound} and the surrogate differs from it by at
most $a$ after conditional maximization. Hence all relevant suprema are
finite. Bisection and a sample at a strictly smaller threshold, exactly
as in \eqref{eq:a-wcac-certified-interval}, produce an enclosing interval
of width at most $\eps/4$ and an algebraic core $(z_*,q_*,v_*)$ whose
recoverable physical scenario loses at most $5\eps/16$. The estimates
apply because at fixed $z$ coefficients in different blocks are independent:
the exceptional and nonexceptional optimizers can be combined.

Recover the exceptional coefficients from this feasible core using the
zonotope argument in the proof of \cref{lem:a-blk-boxlp}. That recovery
works for \emph{any} feasible algebraic core and value, not only for the
box lemma's own maximizer: enumerate exposed vertices of the
$(\kappa+1)$-dimensional image, find a convex combination of at most
$\kappa+2$ of them representing the prescribed image, and lift that
combination to original box coordinates. Every operation uses the one
ordered field of $(z_*,q_*,v_*)$, of polynomial degree and encoding.
Recover each selected nonexceptional cycle profile in its own local
extension of the field of $z_*$, as in \cref{sec:a-weighted-cactus}.
This gives consistent local physical states at the same nomination.

Finally isolate the coordinates of $z_*$ to width $h$ and solve rational
linear feasibility in the intersection of that isolating box with $P$.
It is nonempty because it contains $z_*$. Round all $k\le2m$ original
coefficient coordinates inside their original intervals to width
$\delta$. With $H_1=\sum_{v,j}|H_{vj}|$, choose positive rational widths
such that
\[
 K_bH_1h+\|c\|_1B_0^2k\delta\le\eps/4,
\]
where $K_b$ is from \eqref{eq:a-wcac-nomination-bound}. Polynomial-bit
choices suffice, including zero coefficients of either term. The
all-graph continuity estimate \eqref{eq:a-wcac-joint-continuity} bounds
the remaining loss. Rounding need not retain the selected algebraic
sign part or its circulation: the original rational parameters have
their own unique physical state. The total loss is below $\eps$.
All counts, representations, and precision demands are polynomial for
fixed $d,\kappa_0$. The minimum follows by changing $c$ to $-c$.
\end{proof}

\begin{corollary}[Balanced boxes with fixed support and exceptional rank]
\label{cor:a-wblk-hybrid-box}
Fix $p,\kappa_0$. If $\kappa(G)\le\kappa_0$, nominations range over a
nonempty rational balanced box, and $|\supp c|\le p$, then the value and
rational-scenario guarantees of \cref{thm:a-wblk-hybrid} hold in each
of the three coefficient models.
\end{corollary}
\begin{proof}
Every block rank is at most $\max\{1,\kappa_0\}$. Apply
\cref{thm:a-wcac-faces}; pruning and zero-objective block contraction
cannot increase $\kappa$. Each of polynomially many faces has a fixed
number of affine nomination parameters. Apply
\cref{thm:a-wblk-hybrid} with tolerance $\eps/8$ on each nonempty face.
Take the maximum of the certified lower endpoints and the maximum of
the upper endpoints, and select the scenario associated with the largest
lower endpoint. This loses at most one additional interval width.
The resistance-independent face property and rational disaggregation
in \cref{thm:a-wcac-faces} complete the joint original-network guarantee.
The zero objective is handled by rational feasibility alone.
\end{proof}

\subsection{The remaining several-parameter summation question}

The approximation ingredients have substantial predecessors. Vigneron's
Theorem~6 and Section~3.2 give a multiplicative approximation scheme for
sums of nonnegative algebraic functions of constant description complexity,
using common arrangements and approximate values; Section~2.3 states a
bit-model extension \citep{vigneron2014-geometric-optimization-and-sums-of}.
Its accuracy dependence is polynomial in $1/\eps$. Our scalar lemma
instead uses Lipschitz control near exceptional parameters and geometric
analytic panels to obtain polynomial dependence on accuracy bits. It
supplies rational polynomials in the \emph{original common parameter},
so they can be added without introducing a new variable per block.
No new root-isolation, complex inverse, or interpolation principle is
claimed. Verified enclosure of piecewise analytic functions also has
direct predecessors: Petras supplies such approximation algorithms under
assumptions on the function representation \citep{petras2002-approximation}.

Borcea, B\o gvad, and Shapiro give rational approximants converging
exponentially to a dominant algebraic branch away from an exceptional
locus \citep[Theorems~2--3]{borcea2006-rational-algebraic}.
That statement does not itself give uniform rational approximations
across the exceptional parameters with the present coefficient-bit
and branch-selection guarantees.

Yomdin gives a bound of order $\log^3(1/\eps)$ for the degree-based
complexity of parametric polynomial approximations to planar semialgebraic
sets \citep[Definitions~5.1--5.2 and Theorem~5.6]{yomdin2014-parametrizations}.
This approximates a parametrization of a set, rather than the value of
a function in a prescribed original coordinate. In several variables,
Binyamini and Novikov give polynomial bounds for the number and
semialgebraic complexity of smooth parameterizing charts
\citep[Theorem~1]{binyamini2019-complex-cells}. Such charts parameterize
each set separately; their stated geometric bounds do not by themselves
supply a rational-coefficient bit algorithm for adding many value
functions in their original shared coordinates. These are distinctions
between the inspected statements and the needed interface, not claims
that approximation theory cannot supply it.

More precisely, after the face reduction and
\cref{lem:a-wblk-local}, the unresolved case in the present algorithms is
\begin{equation}
 \max_{z\in P}\sum_{i=1}^m V_i(z),\qquad d\ge2\text{ fixed},
 \label{eq:a-wblk-residual}
\end{equation}
where $m$ is unrestricted, $P\subseteq\R^d$ is a compact rational
polytope, and each $V_i$ is a continuous bounded Lipschitz semialgebraic
block-value function with a polynomial-size bounded-degree rational graph
description and polynomial-bit bounds. A polynomial-size, uniformly
accurate rational piecewise approximation in the shared $z$ coordinates,
constructible in accuracy-bit polynomial time, would suffice: take a
common sign decomposition, add the rational surrogates, optimize in fixed
dimension, and use the global continuity estimate for rational recovery.
An optimization method that bypasses such representations could also
suffice. Neither construction is asserted here for arbitrary $m,d$.

The obstruction is no longer univariate root isolation or the nomination
face reduction; those have been addressed above. A Lipschitz grid in
$d$ dimensions has a power of $1/\eps$ points. Directly retaining one
local value or circulation vector per block in a quantifier-elimination
formula loses the fixed-variable bound. Counts of zeros of exact
univariate sums alone supply neither effective near-cancellation control
nor a several-variable optimizer. These observations explain why those
particular arguments do not prove an accuracy-bit bound. They establish
no impossibility result, no new hardness classification, and no necessity
to solve the exact SRS comparison problem for additive optimization.
The proved higher-block guarantees are the scalar-parameter theorem and
the several-parameter theorem with fixed exceptional rank, together with
the earlier full cactus and fixed-total-rank results.
